\documentclass[12pt,a4paper]{article}

\usepackage[T2A]{fontenc}
\usepackage[utf8]{inputenc}
\usepackage[russian]{babel}
\usepackage{amsmath,amssymb,amsfonts}
\usepackage{booktabs}
\usepackage{longtable}
\usepackage{geometry}
\usepackage{array}
\usepackage{enumitem}

\geometry{margin=2.2cm}

\title{Математическая модель IISVersion5}
\author{Проект анализа интегрального индекса состояния}
\date{\today}

\begin{document}
\maketitle

\section{Назначение документа}
Настоящий документ фиксирует математическую формулу версии \textbf{IISVersion5}, введённой как развитие \textbf{IISVersion4}. Основная цель новой версии состоит не в добавлении новых физиологических блоков, а в уменьшении \textbf{плоскости интерпретации} итоговой шкалы. Версия 5 сохраняет опору на реально доступные EEG и ECG/HRV признаки, но заменяет слишком мягкую итоговую интерпретацию на более контрастную нелинейную схему.

\section{Обозначения}
Используются следующие обозначения:
\begin{itemize}[leftmargin=1.25cm]
    \item $A$ --- EEG-компонент латерализации и асимметрии.
    \item $\Gamma$ --- EEG-компонент активационной нагрузки по gamma-диапазону.
    \item $V$ --- вегетативный компонент по ECG/HRV.
    \item $Q$ --- интегральный физиологический компонент согласованности состояния.
    \item $IIS_{V5}$ --- итоговый интегральный индекс состояния версии 5.
    \item $\varepsilon$ --- малый стабилизирующий параметр.
    \item $\sigma(x)=\dfrac{1}{1+e^{-x}}$ --- логистическая сигмоида.
    \item $\tanh(x)$ --- гиперболический тангенс.
\end{itemize}

Также используется оператор мягкой центрированной нормировки:
\begin{equation}
T(x;c,w)=\tanh\!\left(\frac{x-c}{w}\right),
\end{equation}
где $c$ --- центр калибровки, а $w$ --- характерная ширина диапазона.

\section{Компонент \(A\)}
Компонент $A$ строится как смесь alpha-асимметрии и общей межполушарной асимметрии.

Сначала alpha-асимметрия:
\begin{equation}
A_{\alpha}=T(\alpha_{\mathrm{asym}};\ 0.0006126059,\ 0.0005296437).
\end{equation}

Затем асимметрия суммарной мощности:
\begin{equation}
A_{\mathrm{total}}=
T\!\left(
\frac{P_L-P_R}{P_L+P_R+\varepsilon};\
-0.0298851838,\ 0.0287095439
\right).
\end{equation}

Итоговый компонент:
\begin{equation}
A=\tanh\!\left(1.10\,A_{\alpha}+0.40\,A_{\mathrm{total}}\right).
\end{equation}

\section{Компонент \(\Gamma\)}
Gamma-компонент трактуется как маркер перегрева и активационной нагрузки.

Сначала логарифмируется gamma-мощность:
\begin{equation}
G_{\log}=\ln(1+P_{\gamma}\cdot 10^{12}),
\end{equation}
где множитель $10^{12}$ переводит мощность из $V^2$ в $\mu V^2$.

Далее нормируются две части:
\begin{equation}
G_{\gamma}=T(G_{\log};\ 1.5434688676,\ 0.8496713487),
\end{equation}
\begin{equation}
G_{\gamma/\alpha}=T(\ln(1+\gamma/\alpha);\ 0.0003674267,\ 0.0004343448).
\end{equation}

Итог:
\begin{equation}
\Gamma=\tanh\!\left(0.30\,G_{\gamma}+0.70\,G_{\gamma/\alpha}\right).
\end{equation}

\section{Компонент \(V\)}
Вегетативный компонент строится как комбинация ЧСС и отношения $HF/LF$.

По частоте сердечных сокращений:
\begin{equation}
V_{hr}=T(118.7401955561-HR;\ 0,\ 46.7184548147).
\end{equation}

По вариабельности:
\begin{equation}
V_{hrv}=T(\ln(1+HF/LF);\ 0.9675457417,\ 1.2220370049).
\end{equation}

Итоговый вегетативный компонент:
\begin{equation}
V=\tanh\!\left(0.80\,V_{hr}+0.20\,V_{hrv}\right).
\end{equation}

\section{Новый компонент \(Q\)}
Главное отличие версии 5 от версии 4 состоит в том, что $Q$ больше не является просто мягкой смесью $A$, $V$ и $\Gamma$. Вместо этого вводится \textbf{контрастная физиологическая согласованность}.

\subsection{Линейное ядро}
\begin{equation}
Q_{\mathrm{lin}}=0.50A+0.42V-0.30\Gamma.
\end{equation}

\subsection{Когерентность}
\begin{equation}
Q_{\mathrm{coh}}=A\cdot V-0.50|A-V|-0.18\max(\Gamma,0).
\end{equation}

Здесь:
\begin{itemize}[leftmargin=1.25cm]
    \item произведение $A\cdot V$ усиливает состояние, когда корковая и вегетативная части согласованы;
    \item штраф $|A-V|$ подавляет ситуации, где EEG и HRV указывают в разные стороны;
    \item штраф $\max(\Gamma,0)$ уменьшает оценку при перегретой активации.
\end{itemize}

\subsection{Итоговая формула \(Q\)}
\begin{equation}
Q=\tanh\!\left(1.4\,Q_{\mathrm{lin}}+0.55\,Q_{\mathrm{coh}}\right).
\end{equation}

\section{Итоговая формула \(IISVersion5\)}
\subsection{Линейное ядро}
\begin{equation}
core=0.12A-0.07\Gamma+0.24V+0.57Q.
\end{equation}

\subsection{Центрирование}
\begin{equation}
z=core-0.08.
\end{equation}

\subsection{Контрастный слой}
\begin{equation}
contrast=
2.6z
+0.8z^3
+0.18(Q-V)
-0.06\max(\Gamma,0).
\end{equation}

Здесь:
\begin{itemize}[leftmargin=1.25cm]
    \item член $2.6z$ задаёт базовое растяжение шкалы;
    \item кубический член $0.8z^3$ усиливает края распределения без резких разрывов;
    \item член $0.18(Q-V)$ повышает контраст между интегральным качеством и чисто вегетативной устойчивостью;
    \item тормоз $-0.06\max(\Gamma,0)$ ослабляет состояние при перегретой активации.
\end{itemize}

\subsection{Базовая и контрастная оценки}
\begin{equation}
IIS_{\mathrm{base}}=\sigma(core),
\end{equation}
\begin{equation}
IIS_{\mathrm{contrast}}=0.5+0.5\tanh(contrast).
\end{equation}

\subsection{Финальная смесь}
\begin{equation}
IIS_{V5}=
\operatorname{clip}\!\left(
0.65\,IIS_{\mathrm{base}}+0.35\,IIS_{\mathrm{contrast}},
0,1
\right).
\end{equation}

Таким образом, версия 5 не отказывается от устойчивой логистической интерпретации, а лишь частично смешивает её с более контрастной нелинейной шкалой.

\section{Proxy-режим}
В режиме \texttt{proxy} итоговый компонент $Q$ может быть мягко якорирован по self-report valence/arousal:
\begin{equation}
Q_{\mathrm{anchor}}=
\tanh\!\left(
2.4\left(
0.65\,valence_{\mathrm{norm}}
+0.35(1-arousal_{\mathrm{norm}})
-0.5
\right)
\right).
\end{equation}

Тогда:
\begin{equation}
Q_{\mathrm{proxy}}=0.75Q+0.25Q_{\mathrm{anchor}}.
\end{equation}

После этого в формулу $IIS_{V5}$ подставляется уже $Q_{\mathrm{proxy}}$.

\section{Динамическое продолжение}
В проекте используется причинная динамическая версия, в которой статический $IIS_{V5}(t)$ подаётся на каузальный сглаживающий блок. Отдельная динамическая физиологическая формула пока не вводится; вместо этого применяется следующий принцип:
\begin{itemize}[leftmargin=1.25cm]
    \item сначала считается статический $IIS_{V5}(t)$ на каждом временном окне;
    \item затем строятся каузальные оценки по прошлым окнам;
    \item ухудшение допускается быстрее, восстановление моделируется медленнее.
\end{itemize}

В символическом виде:
\begin{equation}
IIS_{\mathrm{dyn}}(t)=(1-w_t)\,IIS_{V5}(t)+w_t\,S_t,
\end{equation}
где $S_t$ --- каузальная сглаженная часть, вычисленная только по прошлым окнам, а $w_t$ зависит от текущей волатильности траектории.

\section{Интерпретация}
Смысл версии 5 таков:
\begin{enumerate}[leftmargin=1.25cm]
    \item физиологическое ядро по-прежнему полностью строится на реально измеряемых EEG и ECG/HRV признаках;
    \item середина шкалы перестаёт быть излишне плоской;
    \item переходы между режимами становятся визуально и аналитически более различимыми;
    \item модель не вводит недоступные гормональные измерения и не опирается на гипотезу о коэффициенте $K$.
\end{enumerate}

\section{Эмпирический смысл}
На текущих реальных данных версия 5 показала:
\begin{itemize}[leftmargin=1.25cm]
    \item заметное расширение диапазона итогового индекса по сравнению с версией 4;
    \item более контрастное разделение baseline/stress на датасете ds002724;
    \item ограниченные, но всё ещё не окончательные признаки третьего режима;
    \item отсутствие искусственного подтверждения 4 и более устойчивых состояний.
\end{itemize}

Следовательно, \textbf{IISVersion5} следует трактовать как \textbf{контрастную безгормональную интерпретацию} модели ИИС, а не как новый физиологический блок.

\end{document}
